\documentclass[10pt,a4paper]{amsart}
\usepackage[margin=19mm]{geometry}
\usepackage{amsmath,amssymb,mathtools}
\usepackage[scr=rsfs,cal=euler]{mathalfa}
\usepackage{xfrac}
\usepackage[unicode,hidelinks,bookmarks=false]{hyperref}
\newcommand{\ie}{i.e.\ }
\newcommand{\cf}{cf.\ }
\newcommand{\resp}{respectively\ }
\newcommand{\Subsection}{\subsection}
\providecommand{\nobreakdash}{\nobreak\hbox{-}}
\setlength{\emergencystretch}{2em}
\multlinegap0pt
\usepackage[shortlabels]{enumitem}
\usepackage{mathtools}
\usepackage{amssymb}
\newtheorem{proposition}{Proposition}
\newtheorem{lemma}{Lemma}
\title{Asymptotic Universal Koszulity in Galois Cohomology}
\author{Marina Palaisti}
\date{Source: arXiv:2603.27032v3 (CC BY 4.0)}
\begin{document}
\maketitle

\noindent\emph{Source and licence.} Asymptotic Universal Koszulity in Galois Cohomology. Version arXiv:2603.27032v3 (CC BY 4.0), distributed by arXiv under the Creative Commons Attribution 4.0 International licence, http://creativecommons.org/licenses/by/4.0/, as recorded in the arXiv licence metadata for this exact version and shown on the abstract page https://arxiv.org/abs/2603.27032v3. Full licence notice: \texttt{licenses/arxiv-2603.27032-CC-BY-4.0.txt}.\par\smallskip

\noindent\textit{Editorial scope.} Counted unit: the article's proposition prop:galois (source0.tex lines 324--334). The article's complete original proof of it is delivered with it (source0.tex lines 336--352, one \texttt{proof} environment of 17 lines). No mathematics is written, repaired or re-derived: the statement and the proof are the article's own lines, in the article's own notation, and no retained line is substituted or re-flowed.  Retained uncounted as the source-local context the proof needs: lines 147--154; lines 281--302.  Nothing was omitted from any retained span, so the package carries no omission record.  No retained line contains a citation, so the wrapper carries no bibliography and no bibliography entry.  No retained line contains a figure environment, an image-inclusion command or drawing code, so no image is delivered and the package declares no asset dependency.  Editorial formatting only, in addition: an AMS \texttt{amsart} preamble that restates the article's own declarations and macros used by the retained lines, namely the environments \texttt{proposition}, \texttt{lemma} on separate counters, together with the standard packages those lines need.  The retained environments print in this document as Proposition 1, Lemma 1 in order of appearance, whereas the article prints the same items with the numbering of its own full document.  The complete native source of the article as distributed in the arXiv e-print for this exact version is bundled unchanged as \texttt{sources/197321/source0.tex}.
	\begin{proposition}\label{prop:local}
		Let \(A\) be a connected graded $k$-algebra generated in degree \(1\). Assume that for every finite-dimensional subspace \(W \subset A_1\), there exists a quadratic subalgebra \(B_W \subset A\), such that:
		\begin{enumerate}[label=(\alph*)]
			\item \(W \subset (B_W)_1\) and \(\dim_k (B_W)_1 < \infty\);
			\item \(B_W\) is universally Koszul;
		\end{enumerate}
		and assume moreover that the family \((B_W)\) can be chosen to be filtered under inclusion. Then \(A\) is asymptotically universally Koszul.
	\end{proposition}

	\begin{lemma}\label{lem:cohom-colimit}
		Let \(G\) be a pro-\(p\) group and let \((N_i)_{i\in I}\) be a cofinal directed system of closed normal subgroups of \(G\), such that
		\[\bigcap_{i\in I} N_i = 1.\]
		Set \(G_i:=G/N_i\). Assume:
		\begin{enumerate}[label=(\roman*)]
			\item \(H^\bullet(G)\) is generated as an \(\mathbb{F}_p\)-algebra by \(H^1(G)\),
			\item for each \(i\in I\) and each \(n\in\{1,2\}\), the inflation map 
			\[\mathrm{inf}^n_{G_i}\colon H^n(G_i)\longrightarrow H^n(G)\]
			is injective.
		\end{enumerate}
		Then the inflation maps induce:
		\begin{enumerate}[label=(\alph*)]
			\item isomorphisms
			\[\varinjlim_{i\in I} H^n(G_i)\xrightarrow{\sim} H^n(G) \qquad \text{for } n=0,1,2;\]
			\item surjections
			\[ \varinjlim_{i\in I} H^n(G_i)\twoheadrightarrow H^n(G) \qquad \text{for all } n\ge 0,\]
			compatible with cup products.
		\end{enumerate}
		In particular, the induced homomorphism of graded \(\mathbb{F}_p\)-algebras
		\[\varinjlim_{i\in I} H^\bullet(G_i)\longrightarrow H^\bullet(G)\]
		is surjective and is an isomorphism in degrees \(\le 2\).
	\end{lemma}

	\begin{proposition}\label{prop:galois}
		Let \(G\) be a pro-\(p\) group and let \((N_i)_{i\in I}\) be a cofinal directed system of closed normal subgroups of \(G\) with \(\bigcap_i N_i=1\). Set \(G_i:=G/N_i\). Assume:
		\begin{enumerate}[label=(\roman*)]
			\item \(H^\bullet(G)\) is generated as an \(\mathbb{F}_p\)-algebra by \(H^1(G)\),
			\item for each \(i\) and each \(n\in\{1,2\}\), the inflation map
			\[\mathrm{inf}^n_{G_i}\colon H^n(G_i)\longrightarrow H^n(G)\]
			is injective,
			\item for each \(i\), the algebra \(H^\bullet(G_i)\) is a finite-type universally Koszul \(\mathbb{F}_p\)-algebra.
		\end{enumerate}
		Then \(H^\bullet(G)\) is asymptotically universally Koszul.
	\end{proposition}

	\begin{proof}
		For each \(i\), let
		\[W_i:=\mathrm{inf}^1_{G_i}\bigl(H^1(G_i)\bigr)\subseteq H^1(G),\] and let
		\[B_i:=\langle W_i\rangle_{H^\bullet(G)}\]
		denote the graded subalgebra of \(H^\bullet(G)\) generated by \(W_i\). We claim that the inflation map induces an isomorphism of quadratic algebras
		\[H^\bullet(G_i)\xrightarrow{\sim} B_i.\]
		Indeed, since \(H^\bullet(G_i)\) is quadratic, it is generated by \(H^1(G_i)\), so its image is precisely \(B_i\). It remains to identify the quadratic relation space. In \(H^\bullet(G_i)\), the quadratic relations are
		\[R_i=\ker\!\Bigl(H^1(G_i)^{\otimes 2}\xrightarrow{\cup} H^2(G_i)\Bigr).\]
		Since inflation is injective on \(H^1\) and \(H^2\), and is compatible with cup products, the image of \(R_i\) inside \(W_i^{\otimes 2}\) is exactly
		\[\ker\!\Bigl(W_i^{\otimes 2}\xrightarrow{\cup} H^2(G)\Bigr),\]
		which is the quadratic relation space of \(B_i\). Hence $H^\bullet(G_i)\cong B_i$ as quadratic algebras. Therefore each \(B_i\) is a finite-type universally Koszul quadratic subalgebra of \(H^\bullet(G)\).
		
		Next, if \(N_j\subseteq N_i\), then the quotient map \(G_j\twoheadrightarrow G_i\) implies $W_i\subseteq W_j$, hence $B_i\subseteq B_j$. hus the family \((B_i)\) is filtered under inclusion.
		
		Finally, by Lemma~\ref{lem:cohom-colimit}, every class in \(H^1(G)\) comes from some \(H^1(G_i)\), so \[H^1(G)=\bigcup_i W_i.\]
		Since \(H^\bullet(G)\) is generated by \(H^1(G)\), every homogeneous element of \(H^\bullet(G)\) lies in some \(B_i\). Therefore \[H^\bullet(G)=\bigcup_i B_i.\] Hence \(H^\bullet(G)\) is a filtered union of finite-type universally Koszul quadratic subalgebras. By Proposition~\ref{prop:local}, it follows that \(H^\bullet(G)\) is asymptotically universally Koszul.
	\end{proof}

\end{document}
